\section{Experiment Details}
\label{app:exp}
\subsection{Evaluation}
\label{app:eval_normalization}
One complication that arises during evaluation on real-world websites is that multiple elements on a webpage may induce the same effect. For instance, a button might house a text span within it, both of which, when clicked, yield identical results.
To enhance the robustness of our evaluation, we employ heuristics to detect elements equivalent to the ground truth.
We first examine the ancestors of the labeled element to identify potential higher-level elements acceptable for the current action. We employ a straightforward heuristic that locates the nearest clickable element
to the ground truth, including itself. After identifying the top-level acceptable element, we include all its visible descendants that are located within its post-rendering bounding box as acceptable as well. Manual checking on \num{100} instances where the heuristic identifies a top-level element other than the ground truth confirms the validity of the approach.
 For both training and evaluation stages, all acceptable elements are considered positive.

\subsection{Model Implementation Details}

\begin{table}[t]
    \centering
    \small
    \caption{Hyperparameters used in experiments.}
    \begin{tabular}{rp{3cm}p{6cm}}
    \toprule
        \textbf{Method} & \textbf{Model} & \textbf{Hyperparamerters} \\
        \midrule
        Candidate Generation & \texttt{deberta-v3-base} &\texttt{batch\_size:\num{32}, epoch:\num{5}, learning\_rate:\num{3e-5}}\\\addlinespace
        Action Prediction \\
        \hspace{2em}Generation& \texttt{flan-t5-base} &\texttt{batch\_size:\num{32}, epoch:\num{5}, learning\_rate:\num{5e-5}}\\\addlinespace
        \multirow{1}{*}{\hspace{2em}\model}& \texttt{flan-t5-base, flan-t5-large, flan-t5-xl} &\texttt{batch\_size:\num{32}, epoch:\num{5}, learning\_rate:\num{5e-5}}\\
        &\texttt{gpt-3.5-turbo, gpt-4}&\texttt{temperature:\num{0}, \#~demonstrations:\num{3}}\\
    \bottomrule
    \end{tabular}
    \label{tab:hyper}
\end{table}

\noindent\textbf{Candidate Generation.}
We use the Cross-Encoder implementation from Sentence-Transformers~\footnote{\url{https://www.sbert.net/examples/applications/cross-encoder/README.html}} and use DeBERTa as the backbone model. More specifically, we use \texttt{DeBERTa-v3-base}~\footnote{\url{https://huggingface.co/microsoft/deberta-v3-base}} for our experiments.

\noindent\textbf{Action Prediction.}
We use the Seq2Seq model implementation from Transformers~\cite{wolf-etal-2020-transformers}. We experiment with the \texttt{base}~\footnote{\url{https://huggingface.co/google/flan-t5-base}}, \texttt{large}~\footnote{\url{https://huggingface.co/google/flan-t5-large}} and \texttt{xl}~\footnote{\url{https://huggingface.co/google/flan-t5-xl}} versions of Flan-T5~\cite{flant5}.

\noindent\textbf{LLM In-context Learning.}
We use the OpenAI API for in-context learning with LLMs. We experiment with two versions of GPT models: \texttt{gpt-3.5-turbo} and \texttt{gpt-4}. We include three demonstration examples for in-context learning. The complete prompt is shown in Table~\ref{tab:action_prediction_prompt}.

\noindent\textbf{Training Details.}
The \texttt{flan-t5-xl} and \texttt{flan-t5-large} models are trained on servers with 4*A100 80GB cards provided by Ohio Supercomputer Center~\cite{OhioSupercomputerCenter1987}. All other models are trained with single A6000 48GB cards.

Please see Table~\ref{tab:hyper} for all hyperparameters used in our experiments.


% \end{table}
